\subsection{\texorpdfstring{Homomorphisms \(\operatorname{Tor}^{\mathbf{B}}(\mathbf{P},\mathbf{N})\otimes_{\mathbf{A}}\mathbf{Q}\to \operatorname{Tor}^{\mathbf{B}}(\mathbf{P},\mathbf{N}\otimes_{\mathbf{A}}\mathbf{Q})\) and \(\operatorname{Ext}_{\mathbf{B}}(\mathbf{M},\mathbf{N})\otimes_{\mathbf{A}}\mathbf{Q}\rightarrow \operatorname{Ext}_{\mathbf{B}}(\mathbf{M},\mathbf{N}\otimes_{\mathbf{A}}\mathbf{Q})\)}{Homomorphisms relating Tor and Ext to tensor products}}

Let B be a ring, N a (B, A)-bimodule, M a left B-module, P a right B-module, and Q a left A-module.

According to X, p.~62, we have a homomorphism

\[
\gamma_ {1}: \mathrm{H} (\mathrm{L} (\mathrm{P}) \otimes_ {\mathrm{B}} \mathrm{N}) \otimes_ {\mathrm{A}} \mathrm{Q} \rightarrow \mathrm{H} (\mathrm{L} (\mathrm{P}) \otimes_ {\mathrm{B}} \mathrm{N} \otimes_ {\mathrm{A}} \mathrm{Q});
\]

moreover (X, p.~69, prop. 5), we have isomorphisms

\[
\begin{array}{c} \psi_ {1} (N): \operatorname{Tor} ^ {B} (P, N) \otimes_ {A} Q \qquad H (L (P) \otimes_ {B} N) \otimes_ {A} Q, \\ \psi_ {1} (N \otimes_ {A} Q): \operatorname{Tor} ^ {B} (P, N \otimes_ {A} Q) \to H (L (P) \otimes_ {B} N \otimes_ {A} Q). \end{array}
\]

The canonical degree-0 graded homomorphism

\[
\operatorname{Tor} ^ {\mathrm{B}} (\mathrm{P}, \mathrm{N}) \otimes_ {\mathrm{A}} \mathrm{Q} \rightarrow \operatorname{Tor} ^ {\mathrm{B}} (\mathrm{P}, \mathrm{N} \otimes_ {\mathrm{A}} \mathrm{Q})\tag{9}
\]

is defined as the composite \(\psi_{\mathrm{P}}(\mathbf{N}\otimes_{\mathbf{A}}\mathbf{Q})^{-1}\circ \gamma_{1}\circ (\psi_{\mathrm{P}}(\mathbf{N})\otimes 1_{\mathbf{Q}})\) .

Similarly, we deduce from the canonical morphism of complexes

\[
\alpha : \operatorname{Homgr} _ {B} (L (M), N) \otimes_ {A} Q \rightarrow \operatorname{Homgr} _ {B} (L (M), N \otimes_ {A} Q)
\]

a homomorphism H(α); we have the canonical homomorphism (X, p.~62)

\[
\gamma_ {2}: \mathrm{H} (\operatorname{Homgr} _ {\mathrm{B}} (L (M), N)) \otimes_ {A} Q \rightarrow \mathrm{H} (\operatorname{Homgr} _ {\mathrm{B}} (L (M), N) \otimes_ {A} Q),
\]

and isomorphisms (X, p.~88, prop. 5)

\[
\begin{array}{c} \varphi_ {\mathrm{M}} (\mathrm{N}): \mathrm{H} \big (\operatorname{Homgr} _ {\mathrm{B}} (\mathrm{L} (\mathrm{M}), \mathrm{N}) \big) \otimes_ {\mathrm{A}} \mathrm{Q} \to \operatorname{Ext} _ {\mathrm{B}} (\mathrm{M}, \mathrm{N}) \otimes_ {\mathrm{A}} \mathrm{Q}, \\ \varphi_ {\mathrm{M}} (\mathrm{N} \otimes_ {\mathrm{A}} \mathrm{Q}): \mathrm{H} \big (\operatorname{Homgr} _ {\mathrm{B}} (\mathrm{L} (\mathrm{M}), \mathrm{N} \otimes_ {\mathrm{A}} \mathrm{Q}) \big) \to \operatorname{Ext} _ {\mathrm{B}} (\mathrm{M}, \mathrm{N} \otimes_ {\mathrm{A}} \mathrm{Q}). \end{array}
\]

The canonical degree-0 graded homomorphism

\[
\operatorname{Ext} _ {\mathbf {B}} (\mathbf {M}, \mathbf {N}) \otimes_ {\mathbf {A}} \mathbf {Q} \rightarrow \operatorname{Ext} _ {\mathbf {B}} (\mathbf {M}, \mathbf {N} \otimes_ {\mathbf {A}} \mathbf {Q})\tag{10}
\]

is defined as the composite

\[
\varphi_ {\mathrm{M}} (\mathrm{N} \otimes_ {\mathrm{A}} \mathrm{Q}) \circ \mathrm{H} (\alpha) \circ \gamma_ {2} \circ \left(\varphi_ {\mathrm{M}} (\mathrm{N}) \otimes 1 _ {\mathrm{Q}}\right) ^ {- 1}.
\]

PROPOSITION 7. --- a) If the A-module Q is flat, the homomorphism (9) is bijective. b) If the A-module Q is finitely generated projective, the homomorphism (10) is bijective.

\begin{enumerate}
\def\labelenumi{\alph{enumi})}
\setcounter{enumi}{2}
\item
  If the A-module Q is flat, the ring B is Noetherian, and the B-module M is finitely generated, then the homomorphism (10) is bijective.
\item
  If Q is flat, \(\gamma_{1}\) is bijective (X, p.~66, cor. 2).
\item
  If Q is finitely generated projective, it is flat, hence \(\gamma_{2}\) is bijective, and moreover \(\alpha\) is bijective (II, p.~75, prop. 2, a)).
\item
  Under the hypotheses of c), \(\gamma_{2}\) is bijective since Q is flat. Moreover (X, p.~53, prop. 6), there exists a resolution L of M such that each \(\mathbf{L}_{n}\) be finitely generated free; let \(u: \mathbf{L}(\mathbf{M}) \to \mathbf{L}\) be a homotopy equivalence (X, p.~49, cor. to prop. 3); in the commutative diagram,
\end{enumerate}

\[
\begin{array}{c c c} \operatorname{Homgr} _ {\mathrm{B}} (\mathrm{L} (\mathrm{M}), \mathrm{N}) \otimes_ {\mathrm{A}} \mathrm{Q} & \xrightarrow {\alpha} & \operatorname{Homgr} _ {\mathrm{B}} (\mathrm{L} (\mathrm{M}), \mathrm{N} \otimes_ {\mathrm{A}} \mathrm{Q}) \\ \uparrow & & \uparrow \\ \operatorname{Homgr} _ {\mathrm{B}} (\mathrm{L}, \mathrm{N}) \otimes_ {\mathrm{A}} \mathrm{Q} & \xrightarrow {\bar {\alpha}} & \operatorname{Homgr} _ {\mathrm{B}} (\mathrm{L}, \mathrm{N} \otimes_ {\mathrm{A}} \mathrm{Q}). \end{array}
\]

the vertical arrows deduced from \(u\) are homotopy equivalences (X, p.~64, prop. 3 and p.~83, prop. 3) and \(\overline{\alpha}\) is bijective (II, p.~75, prop. 2 (ii)); hence H(α) is bijective, and the homomorphism (10) is bijective.

